% il existe plusieures classes de documents
% pour des documents plus longs, vous pouvez utiliser
% book ou report
\documentclass[a4paper]{article} 
\usepackage{graphicx}
\usepackage{amsmath}
\usepackage{amssymb}
\usepackage{biblatex} %Imports biblatex package
\usepackage{subfig}
\usepackage{parskip}
\usepackage{multirow}
\usepackage{multirow}
\usepackage[most]{tcolorbox}
\usepackage{xcolor}
\addbibresource{stg3a.bib} %Import the bibliography file
\bibliography{stg3a}
\usepackage{hyperref}

% vous pouvez changer les paramètres : voici les options dispoinibles :
% - a4paper
% - fancysections
% - notitlepage ou titlepage
% - onside ou twoside selon si vous voulez l'imprimer en recto-verso ou en recto
% - sectionmark
% - chaptermark (pour les 
% - pagenumber
% - enmanquedinspiration
% en cas de doutes, pas de doutes, la documentation est sur :
%  https://gitlab.binets.fr/typographix/polytechnique/-/blob/master/source/polytechnique.pdf



\usepackage[a4paper,  fancysections,  titlepage]{polytechnique}
\usepackage[french]{babel}
\usepackage[T1]{fontenc}
\usepackage{blindtext}
\usepackage[hidelinks]{hyperref}

\usepackage{listings}
\usepackage{tikz}
\usepackage{xcolor}
\lstset{
  language=Python,
  basicstyle=\ttfamily\small,
  breaklines=true,
  showstringspaces = false,
  keywordstyle=\color{blue},
  stringstyle=\color{red},
  commentstyle=\color{brown},
  morecomment=[l][\color{magenta}]{\#},
  inputencoding=utf8,
  extendedchars=true,
  literate={à}{{\`a}}1 {ç}{{\c{c}}}1 {é}{{\'e}}1 {è}{{\`e}}1 {ê}{{\^e}}1 {ë}{{\"e}}1 {î}{{\^i}}1 {ï}{{\"i}}1 {ô}{{\^o}}1 {ö}{{\"o}}1 {ù}{{\`u}}1 {û}{{\^u}}1 {ü}{{\"u}}1 {ÿ}{{\"y}}1 {À}{{\`A}}1 {Ç}{{\c{C}}}1 {É}{{\'E}}1 {È}{{\`E}}1 {Ê}{{\^E}}1 {Ë}{{\"E}}1 {Î}{{\^I}}1 {Ï}{{\"I}}1 {Ô}{{\^O}}1 {Ö}{{\"O}}1 {Ù}{{\`U}}1 {Û}{{\^U}}1 {Ü}{{\"U}}1 {Ÿ}{{\"Y}}1 {\_}{\textunderscore}{1}
}

% nous avons défini deux commandes :
\newcommand{\code}[1]{%
    \mbox{\ttfamily%
        \detokenize{#1}%
    }%
}

\newcommand{\resultat}[1]{%
    \quad \rightsquigarrow \quad #1%
}

\DeclareRobustCommand{\ReLU}{\mathrel{\text{%
    \begin{tikzpicture}[baseline=(current bounding box.south)]
        \draw[thick] (0,0) -- (0.25,0) -- (0.4,0.25);
    \end{tikzpicture}}}\!\!}

\newcommand{\E}{\mathbb{E}}
\DeclareMathOperator{\Var}{Var}
\newcommand{\indic}{\mathbb{1}}
\newcommand{\dd}{\,d}

\author{Janis AIAD}
\date{\today}
\title{Analyse de la dynamique d'apprentissage des réseaux de neurones profonds à travers le prisme du Noyau de Tangente Neuronale}
\subtitle{Contributions à la compréhension des corrections à largeur finie et des régimes d'apprentissage}
% pour changer de logo, ajoutez l'image dans un fomat PDF
% ou image en la glissant à droite et remplacez typographix
% par le nom de l'image, si vous ne voulez pas de logo, 
% supprimze la ligne. 

\begin{figure*}[b]
    \includegraphics[width=0.4\linewidth]{logoaividence.jpg}
    \label{fig:enter-label}
\end{figure}


\fontsize{40pt}{48pt}

\begin{document}
	\maketitle 
	
	\tableofcontents
\newpage


\addcontentsline{toc}{section}{Résumé}
\section*{Résumé}
Ce rapport présente les résultats d'un stage de recherche axé sur l'étude de la dynamique d'apprentissage des réseaux de neurones profonds. Nous explorons les mécanismes sous-jacents à travers le formalisme du Noyau de Tangente Neuronale (NTK), en portant une attention particulière aux corrections à largeur finie. Nos travaux s'articulent autour de la décomposition du NTK via l'entraînement de Sobolev, l'analyse de son espace à noyau reproduisant (RKHS) et l'étude des valeurs propres extrêmes. Nous présentons des bornes théoriques pour la plus petite valeur propre du NTK en nous appuyant sur des techniques de matrices aléatoires (RMT). Ces résultats sont ensuite appliqués à l'analyse de différentes architectures, notamment les réseaux de neurones à propagation avant (FCNN), les réseaux à mémoire matricielle (MMNN) et les transformeurs. Nous discutons également des implications de nos découvertes pour les applications en calcul scientifique (SciML) et proposons des pistes pour de futures recherches.

\newpage
\addcontentsline{toc}{section}{Executive summary}
\section*{Executive summary}
(Executive summary placeholder)

\newpage
\addcontentsline{toc}{section}{Introduction}
\section*{Introduction}
L'étude de la dynamique d'apprentissage dans les réseaux de neurones profonds constitue un domaine de recherche fondamental. Le régime du Noyau de Tangente Neuronale (NTK) a fourni un cadre puissant pour analyser le comportement des réseaux infiniment larges, où la dynamique d'entraînement se linéarise. Cependant, les réseaux pratiques ont une largeur finie, ce qui introduit des corrections complexes et un comportement d'apprentissage des caractéristiques (feature learning). Ce rapport vise à approfondir notre compréhension de ces phénomènes.

Nous commençons par une analyse théorique du NTK, en présentant sa décomposition basée sur l'entraînement de Sobolev et en caractérisant son RKHS. Une partie centrale de nos travaux est consacrée à l'étude du spectre du NTK, où nous cherchons à borner ses plus petites valeurs propres à l'aide d'outils issus de la théorie des matrices aléatoires. Ces résultats théoriques sont ensuite mis en perspective à travers l'analyse de plusieurs architectures : FCNN, MMNN et transformeurs. Pour les FCNN, nous étudions les corrections à largeur finie et le scaling des paramètres. Pour les MMNN et les transformeurs, nous discutons des défis spécifiques liés à l'optimisation et des intuitions sur leurs performances.

Enfin, nous explorons les applications de nos travaux, notamment en SciML, et discutons de l'alignement entre les hypothèses théoriques et les observations empiriques. Nous abordons également des perspectives inspirées de la physique statistique pour interpréter la convergence des modèles vers l'équilibre.


\newpage
\addcontentsline{toc}{section}{Notations}
\section*{Notations}
Les notations se veulent claires et, si possible, colorées, afin de s'adresser au plus grand nombre.
$\ReLU$ : La fonction d'activation ReLU (Rectified Linear Unit), définie par $\ReLU(x) = \max(0, x)$.
\newline
$\E[\cdot]$ : Espérance mathématique.
\newline
$\Var(\cdot)$ : Variance d'une variable aléatoire.
\newline
$\indic\{\cdot\}$ : Fonction indicatrice, valant 1 si la condition est vraie, 0 sinon.
\newline
$\dd x$ : Élément différentiel utilisé dans les intégrales.

\newpage

\section{Cadre Théorique du Noyau de Tangente Neuronale}

\subsection{Entraînement de Sobolev et Décomposition du NTK}
(Présentation des théorèmes sur la décomposition du NTK via le Sobolev training. Description du RKHS pour le noyau de Sobolev et le noyau de Laplace.)

\subsection{Analyse Spectrale du NTK et Bornes sur les Valeurs Propres}
(Discussion sur la plus grande valeur propre du NTK. Présentation des théorèmes de Terjek et du plan de preuve pour obtenir une borne sur la plus petite valeur propre en utilisant la RMT.)

\section{Corrections à Largeur Finie et Régimes d'Apprentissage}

\subsection{Développement Théorique pour les FCNN}
(Présentation des théorèmes et formules pour les corrections à largeur finie dans les FCNN. Analyse des lois de scaling du reste et discussion sur les termes d'erreur, en lien avec les ensembles de matrices aléatoires QUE & GOE.)

\subsection{Approches Alternatives et Concentration}
(Explication de l'approche alternative pour les FCNN via les réseaux paraboliques de Terjek, permettant de se placer en régime de feature learning tout en conservant des propriétés du NTK. Discussion sur les bornes de concentration.)

\section{Analyse d'Architectures Spécifiques}

\subsection{Réseaux à Mémoire Matricielle (MMNN) et Transformeurs}
(Description des techniques de préconditionnement du NTK pour ces architectures. Discussion sur les régimes "feature" et "kernel", et sur les défis liés à l'optimisation. Présentation d'une conjecture sur la stabilité autour des minima globaux.)

\subsection{Intuitions sur les Performances}
(Discussion finale basée sur les résultats théoriques et expérimentaux pour expliquer les performances observées pour les MMNN, les transformeurs et les ResNets.)

\section{Résultats Expérimentaux et Applications}

\subsection{Validation des Lois de Scaling et des Corrections}
(Présentation des expériences numériques validant les lois de scaling pour les corrections à largeur finie dans les FCNN.)

\subsection{Analyse Hessienne et Dynamique}
(Présentation des expériences sur l'approche hessienne du NTK et l'étude des termes de type Lyapunov.)

\subsection{Applications en SciML et Discussion}
(Discussion sur l'application des résultats en Scientific Machine Learning. Vérification des hypothèses théoriques sur les grands modèles et le scaling des paramètres. Lien avec la physique statistique et la convergence vers l'équilibre thermique.)


\newpage


\section*{Annexes}

\begin{figure}[h]
    \centering
    \includegraphics[width=0.8\linewidth]{001504621_896x598_c.png}
    \caption{Aperçu de l'interface \href{https://github.com/AI-vidence/antakia}{mise à disposition en source ouverte} pour l'extraction de connaissances}
    \label{fig:enter-label}
\end{figure}


\newpage


\appendix
\section{Expériences numériques}
(paragraphe)


\printbibliography


	
\end{document}